\documentclass[12pt]{article}
\usepackage[T2A]{fontenc}
\usepackage[utf8]{inputenc}
\usepackage[ukrainian]{babel}
\usepackage{amsmath,amssymb}
\usepackage{array}

\title{Практична робота: Математичні формули}
\author{Миколаєнко Дмитро}
\date{\today}

\begin{document}
\maketitle

\section*{Приклади математичних формул}

Вбудовані формули у тексті: $a^2+b^2=c^2$, $D=b^2-4ac$, $\sqrt{16}=4$.

\begin{equation}
E=mc^2
\end{equation}

\begin{align}
a+b &= c\\
x^2+y^2 &= z^2
\end{align}

Приклади дробів та коренів:
\[
\frac{a+b}{c+d}, \qquad \sqrt{x+1}, \qquad \sqrt[3]{8}=2
\]

\section*{Квадратне рівняння}



Запис $D$ --- це вбудована формула в тексті (дискримінант).

\begin{equation}
ax^2+bx+c=0,\quad a\neq0
\end{equation}

\section*{Проста таблиця 3x3}

\begin{table}[h]
\centering
\begin{tabular}{|c|c|c|}
\hline
1 & 2 & 3\\
\hline
4 & 5 & 6\\
\hline
7 & 8 & 9\\
\hline
\end{tabular}
\caption{Проста таблиця}
\end{table}

\section*{Таблиця з математичними виразами}

\begin{table}[h]
\centering
\renewcommand{\arraystretch}{1.3}
\begin{tabular}{|c|c|}
\hline
Вираз & Результат\\
\hline
$a^2$, $a=3$ & 9\\
\hline
$\sqrt{25}$ & 5\\
\hline
$\frac{12}{4}$ & 3\\
\hline
\multicolumn{2}{|c|}{Математичні обчислення}\\
\hline
\end{tabular}
\caption{Таблиця математичних виразів}
\end{table}

\section*{Розв'язання квадратного рівняння}


У цьому блоці показано покрокове розв’язання рівняння.

\begin{align*}
x^2-5x+6 &= 0\\
D &= (-5)^2-4\cdot1\cdot6\\
D &= 25-24\\
D &= 1\\
x_{1,2} &= \frac{5\pm\sqrt{1}}{2}\\
x_1 &= 3\\
x_2 &= 2
\end{align*}

\end{document}
